\documentclass[11pt,a4paper]{article}
\usepackage[T1]{fontenc}
\usepackage{lmodern,amsmath,amssymb,amsthm,mathtools,booktabs,microtype}
\usepackage[margin=27mm]{geometry}
\usepackage[colorlinks=true,linkcolor=blue,citecolor=blue,urlcolor=blue]{hyperref}
\setlength{\emergencystretch}{2em}
\newtheorem{theorem}{Theorem}[section]
\newtheorem{lemma}[theorem]{Lemma}
\newtheorem{proposition}[theorem]{Proposition}
\newtheorem{corollary}[theorem]{Corollary}
\newcommand{\Q}{\mathbb Q}\newcommand{\Z}{\mathbb Z}
\DeclareMathOperator{\Tr}{Tr}\DeclareMathOperator{\Nm}{N}
\title{Fixed trace obstructions for native elliptic detectors}
\author{}\date{16 November 2026}
\begin{document}\maketitle
\begin{abstract}
For the literal general elliptic detector source, the integer input
$X=4$ prevents the two normalized reciprocal traces from both equalling
one, for every guarded auxiliary tuple, every slope, and every odd
scalar exponent. The obstruction persists after any finite sequence
of a class of native normalizations that preserve the detector output.
It uses the prime $252877$, inert in $\Q(\sqrt2)$, at which the scalar
has valuation one. More generally, no finite set of fixed rational
trace targets suffices for all integer inputs in this source.
For either conditional six-contact source, a real scale bound excludes simultaneous unit traces at every prepared slope once the source shift is at least nineteen.
The obstructions concern prescribed trace values; it does not exclude
the adaptive norm equations with varying first coordinates or give a
lower bound for defining-polynomial variables.
\end{abstract}

\section{The original scalar and contact blocks}

Set $K=\Q(\sqrt2)$ and use the literal general source with unsquared
comparison scale in~\cite{Unit}. Its scalar is
\begin{equation}\label{eq:scalar}
\begin{gathered}
 \beta=216X^4+450459,\qquad \ell=76249,\\
 r=\beta(\beta-1)(\beta-2401)(32\beta+343)(2\beta-1),\\
 C_L=17(\ell^{(L+1)/2}+1),\qquad
 \rho=\frac{102^2\ell^Lr}{(1+C_Lr)^2}\quad(L\ge3\text{ odd}).
\end{gathered}
\end{equation}
The fixed exceptional set is
\[
 S=\{p\text{ prime}:p\le76176\}
                  \cup\{1875611,3828397,3982747\}.
\]
In particular $\ell\notin S$. For the rational source height trace $T$
and comparison slope $\lambda$, put
\begin{equation}\label{eq:weights}
\begin{gathered}
 a_0=1,\quad a_1=1+\rho T^2,\quad t_i=a_i\lambda,\quad
 \gamma_i=-i-1,\\
 \Theta=\{0,-1/2,-1,1/2399,-32/407\},\quad
 \Theta_i=(4i+2+\Theta)\cup(4i+4+\Theta),\\
 R_i=\frac{\prod_{\theta\in\Theta_i}(t_i-\theta)}
                  {(t_i-\gamma_i)^{10}},\qquad w_i=\rho R_i
                  \qquad(i=0,1).
\end{gathered}
\end{equation}
Every displayed native denominator and every $w_i$ is guarded nonzero.
The argument below permits arbitrary rational $T,\lambda$ consistent
with these guards and arbitrary guarded native field tests. It
therefore applies, in particular, to all actual elliptic point tuples.

Each detector has the native form
\begin{equation}\label{eq:native}
 W_i=\prod_k(1+w_i h_{i,k}^2),\quad
 V_i=\Tr_{K/\Q}(W_i^{-1}),\quad
 J_i=\Nm_{K/\Q}(W_i^{-1}),\quad b_i=4w_i/V_i^2.
\end{equation}
This is the exact native factorization in~\cite{Adaptive}; the functions
$h_{i,k}\in K$ are not chosen independently in the elliptic source.
For the first normalization, retain the guard $1+w_i\ne0$ and set
\begin{equation}\label{eq:first}
 \omega_i=\frac{w_i}{(1+w_i)^2},\qquad
 \widetilde W_i=(1+w_i)W_i,\qquad
 \widetilde V_i=\frac{V_i}{1+w_i}.
\end{equation}
It preserves $b_i$ and has a native factorization, using the added test
$1+w_i$ and the rescaled original tests $(1+w_i)h_{i,k}$.
Every odd valuation of $\omega_i$ is positive~\cite{Adaptive}.

\section{Iterated native normalizations}

The following operations allow the multiplier to depend on the retained
point data. They introduce no free rational parameter when their
expressions are fixed in advance.

\begin{proposition}\label{prop:iteration}
Suppose $W=\prod_k(1+\omega h_k^2)$ and every odd rational valuation
of $\omega$ is positive. Let $c$ be a rational number, or a specified
native scalar rational function evaluated at a guarded tuple. Define
\begin{equation}\label{eq:iteration}
 G=1+\omega c^2,\qquad
 \omega'=\omega/G^2,\qquad W'=GW,\qquad
 V'=V/G,\qquad J'=J/G^2,
\end{equation}
retaining the guard $G\ne0$. Then $4\omega'/(V')^2=4\omega/V^2$,
and $W'$ has native tests
\[
                         h'_0=cG,\qquad h'_k=Gh_k.
\]
Every odd valuation of $\omega'$ is positive.

At a prepared point-algebra specialization, it is sufficient that
$c$ be defined at every character. It may vanish. The element $G$
is then strictly positive at every character and is a whole-algebra
unit. Only the denominator of $c$, in addition to the inherited
and displayed guards, needs a nonzero guard.
\end{proposition}
\begin{proof}
The added factor is $1+\omega'(cG)^2=G$, and every other factor
is unchanged. The trace, norm, and detector identities follow by
scalar multiplication.

Let $n=v_p(\omega)$ be odd. It is positive by hypothesis. If $c=0$
then $G=1$. Otherwise put $s=v_p(c)$. Since $n+2s$ is odd, the
summands in $G$ have distinct valuations and
\[
 v_p(\omega')=n-2\min\{0,n+2s\}>0.
\]
An even valuation of $\omega$ remains even after subtracting
$2v_p(G)$, so these are all odd cases.

At a prepared specialization, $\omega$ is positive at every geometric
point-algebra character, and all those characters are real~\cite{Adaptive}.
Thus $1+\omega c^2\ge1$ at every character where $c$ is defined.
Nonvanishing at every character gives the unit assertion.
\end{proof}

We call a detector obtained from~\eqref{eq:first} by any finite sequence
of~\eqref{eq:iteration} a \emph{final normalized detector}. Write its
product, reciprocal trace, and weight as $\widehat W_i,\widehat V_i$
and $\widehat w_i$. The choices of the functions $c$ and the lengths
of the sequences may differ between the two blocks.

\begin{lemma}\label{lem:inert}
Let $p\notin S$ be inert in $K$. If a native weight $\omega$ has
odd valuation at $p$, then the reciprocal trace is $p$-integral and
reduces to zero or two modulo $p$.
\end{lemma}
\begin{proof}
The extension at $p$ is unramified. For a nonzero test $h$, the two
summands in $1+\omega h^2$ have distinct valuations, one zero and
one odd. Its reciprocal is integral and reduces to one if
$v_p(\omega)+2v(h)>0$, and to zero otherwise. A zero test gives
the factor one. The reciprocal product consequently reduces to
zero or one. These residues are rational, so their field traces
are zero or two. The proof includes negative odd weights.
\end{proof}

\section{A concrete obstruction at an actual integer input}

\begin{theorem}\label{thm:four}
For the scalar and contacts~\eqref{eq:scalar}--\eqref{eq:weights},
let $X$ be any integer satisfying
\begin{equation}\label{eq:progression}
                  X\equiv4\pmod{63946777129}.
\end{equation}
For every odd $L\ge3$, every guarded rational slope and auxiliary tuple,
and every pair of final normalized detectors, at least one of
$\widehat V_0,\widehat V_1$ is different from one. Consequently
\begin{equation}\label{eq:impossible}
                   \Nm(\widehat W_i-1)=1\qquad(i=0,1)
\end{equation}
cannot hold simultaneously. The assertion includes $X=4$ itself.
\end{theorem}
\begin{proof}
Put $p=252877$. Trial division through $\lfloor\sqrt p\rfloor=502$
verifies primality. One has $p\equiv5\pmod8$, so $p$ is inert in $K$;
also $p\notin S\cup\{\ell\}$. At $X=4$,
\[
 \beta=505755,\quad \beta-1=2p,
\]
and the five scalar factors are
\[
                505755,\ 505754,\ 503354,\ 16184503,\ 1011509.
\]
Their residues modulo $p$ are respectively
\[
                         1,\ 0,\ -2400,\ 375,\ 1.
\]
Thus $v_p(r)=1$. All coefficients are integers, so the same exact
valuation and unit assertions hold throughout~\eqref{eq:progression},
since $p^2=63946777129$. The factors $102,\ell$ are $p$-units and
$1+C_Lr\equiv1\pmod p$ for every odd $L$. Hence
\begin{equation}\label{eq:rhov}
                              v_p(\rho)=1.
\end{equation}

The contact and pole constants in each block are nonzero modulo $p$.
Their residue sets, in the order of~\eqref{eq:weights} followed by the
pole, are
\begin{align*}
\mathcal C_0={}&\{2,126440,1,97084,85744,4,126442,3,97086,85746,252876\},\\
\mathcal C_1={}&\{6,126444,5,97088,85748,8,126446,7,97090,85750,252875\}.
\end{align*}
All denominators are $p$-units and these two sets are disjoint.

We verify the unit-weight alternative directly. If $T=0$, then
$a_1=1$. Otherwise $v_p(\rho T^2)$ is odd, so either
$a_1\equiv1\pmod p$ or $v_p(a_1)<0$. If $v_p(\lambda)\ne0$,
then $R_0$ is a unit: for positive argument valuation all constant
terms are units, and for negative argument valuation the numerator
and denominator have equal degree ten. If $\lambda$ is a unit
and $R_0$ is not, its residue belongs to $\mathcal C_0$. Either
$a_1\lambda$ has negative valuation, or it has the same residue as
$\lambda$ and avoids $\mathcal C_1$. In either case $R_1$ is a unit.
Thus at least one $w_i=\rho R_i$ has valuation one.

Normalization~\eqref{eq:first} and Proposition~\ref{prop:iteration}
preserve odd parity and make the final weight valuation positive.
Lemma~\ref{lem:inert} gives $\widehat V_i\equiv0$ or $2\pmod p$
for this block. It cannot equal one.

Finally, at every guarded tuple,
\[
 \Nm(\widehat W-1)-1
      =\Nm(\widehat W)-\Tr(\widehat W)
      =\Nm(\widehat W)(1-\widehat V).
\]
Since $\Nm(\widehat W)\ne0$, equation~\eqref{eq:impossible} is
exactly the forbidden unit-trace condition.
\end{proof}

The proof does not assume that arbitrary native tests can be realized
by elliptic points. It proves impossibility for every guarded tuple,
and hence for the actual source tuples. It uses neither a height
estimate nor any unproved rational-point realization assertion.

\section{Finite fixed trace targets}

\begin{lemma}\label{lem:primes}
There are infinitely many primes $p$ inert in $K$ for which at least
one of $\beta(X)$, $\beta(X)-1$, $\beta(X)-2401$ has a simple root
modulo $p$. These primes can avoid any prescribed finite set.
\end{lemma}
\begin{proof}
Apply the Chebotarev lemma in~\cite[Section 4]{Finite} with both
square classes equal to $2$. That lemma supplies a simple root of
one of the three displayed quartics and nonsquareness of $2$ modulo
$p$. Its proof uses the distinct imaginary quadratic subfields of the
three quartic root fields, so it applies to the real class $2$.
\end{proof}

\begin{theorem}\label{thm:finite}
Fix a finite set $\mathcal T\subset\Q$. There is an arithmetic
progression of integer inputs, independent of $L$, for which no
prepared guarded tuple of the original scalar source has both
final normalized traces in $\mathcal T$. This remains true for any
fixed finite collection of normalization schemes and trace targets.
\end{theorem}
\begin{proof}
At prepared characters the original native factors are at least one,
while the first normalization multiplies them by $1+w_i>1$.
Subsequent multipliers are at least one. Therefore
\[
                         0<\widehat V_i<2.
\]
It suffices to consider the finite subset
$\mathcal T\cap(0,2)$; if it is empty, there is nothing to prove.

Choose a prime from Lemma~\ref{lem:primes} outside $S\cup\{\ell\}$, outside all
fixed contact denominators and collisions, and outside the numerator
and denominator supports of every $k$ and $k-2$, for
$k\in\mathcal T\cap(0,2)$. Thus no such $k$ reduces to zero or two.
If $\beta(x_0)-\theta$ has a simple root modulo $p$, exactly one
lift of $x_0$ modulo $p^2$ is still a root modulo $p^2$. Choose any
other lift $x_1$. For all $X\equiv x_1\pmod{p^2}$,
$v_p(\beta(X)-\theta)=1$ and the other scalar contacts are units,
after excluding their fixed collisions. Hence $v_p(\rho)=1$
uniformly in the odd exponent $L$.

The unit-weight argument from Theorem~\ref{thm:four} and the inert
residue lemma force at least one final trace to reduce to zero or two.
It cannot belong to $\mathcal T$. For finitely many prescribed
collections, take their finite union and exclude all their fixed
prime exceptions when choosing $p$.
\end{proof}

This theorem rules out a uniform realization argument that replaces
both actual detector traces by values from a finite fixed rational
set. It does not rule out rational points with varying traces. It
also does not apply automatically to a different scalar input map.
In particular the norm-compatible scalar of~\cite{NormScalar} has
different arithmetic and is not covered by the present theorem.

\section{A separate norm restriction on normalization orbits}

\begin{proposition}\label{prop:normorbit}
Let $\omega>0$ and $F=\Q(\sqrt{-\omega})$. In any finite sequence
of~\eqref{eq:iteration}, the product of the multipliers $G$ is a
norm from $F$. Consequently a final trace equal to one requires
the initial reciprocal trace to be a norm from $F$.
\end{proposition}
\begin{proof}
Every new weight differs from $\omega$ by a rational square, so its
quadratic norm field remains $F$. Each multiplier satisfies
$G=\Nm_F(1+c\sqrt{-\omega_{\rm current}})$. Multiply these
identities and use $V_{\rm final}=V_{\rm initial}/\prod G$.
\end{proof}

The restriction is not automatic from the native factorization, even
when all scalar preparation conditions are imposed. For example, set
\[
 s=\frac{100}{1001},\qquad w=s^2,\qquad
 h=\frac{1001}{1000}\sqrt2.
\]
Then
\[
 W=1+wh^2=\frac{51}{50},\qquad
 b=wW^2=\left(\frac{102}{1001}\right)^2\in(0,17).
\]
The input is a rational square and has positive valuation two at each
of $2,3,17$. After the first normalization,
\[
 \widetilde V=\frac{100200100}{51612051},\qquad
 v_3(\widetilde V)=-1.
\]
The norm field is $\Q(i)$; a norm from this field has even valuation
at its inert prime $3$. Thus no sequence of these normalizations
gives final trace one.

Arbitrarily large positive valuations at the three preparation primes
are obtained by taking, for any integer $M\ge2$,
\[
 s=\frac{102^M}{10\cdot102^M+1},\qquad
 h=\frac{\sqrt2}{10s}.
\]
Again $W=51/50$, and $b=(51s/50)^2$ has valuations
$2M-2$ at $2$ and $2M+2$ at both $3$ and $17$.
Since $1+s^2$ is a $3$-adic unit, the normalized trace still has
valuation $-1$ at $3$. These examples are native products with
prepared scalar inputs; they are not asserted to arise from actual
elliptic detector tuples.

\section{The conditional positive outer interval}
\label{sec:unit-trace-positive-interval}

Put
\[
 F(t)=\frac{t(t+1)(2t+1)(3t+1)(4t+1)(5t+2)}{120},\qquad
 w(t)=\frac{F(t)}{4(t-1)^8},\qquad
 \alpha(t)=\frac{(t-1)^4}{F(t)}.
\]
For the conditional six-contact sources in~\cite{Conditional}, fix an actual point tuple. Then the
two real detector products have the form
\[
 W_\pm(t)=\prod_{k=1}^{r}(1+\alpha(t)c_{\pm,k}),
 \qquad c_{\pm,k}\ge0,
\]
where the constants are the squares of the corresponding real test
values; the native source has $r=6$. Define
\[
 V^*(t)=\frac{W_+(t)^{-1}+W_-(t)^{-1}}{1+w(t)}.
\]

\begin{theorem}\label{thm:unit-trace-positive-unique}
For every pair of finite arrays $c_{\pm,k}\ge0$, the equation
$V^*(t)=1$ has exactly one real solution in $t>1$.
Consequently, if $a>1$ and $t>1$, then
\[
                         V^*(t)=V^*(at)=1
\]
is impossible.
\end{theorem}

\begin{proof}
Write $\delta=t-1$ and $f=F'/F$. Since every contact is
nonpositive,
\[
                 0<f<\frac6\delta,\qquad
 \frac{w'}w=f-\frac8\delta<0,\qquad
 \frac{\alpha'}\alpha=\frac4\delta-f.
\]
The polynomial identity
\[
 4F-\delta F'=-\frac{H(t)}{60},\qquad
 H(t)=120t^6-211t^5-745t^4-610t^3-205t^2-28t-1
\]
shows that $\alpha'>0$ implies $t<19/5$. Indeed,
\begin{align*}
 625H(u+19/5)={}&75000u^6+1578125u^5+13273750u^4\\
 &+55806500u^3+117399225u^2+99363501u+1402506
\end{align*}
has positive coefficients. As $w$ is decreasing, the implication
$\alpha'>0$ therefore gives
\begin{equation}\label{eq:unit-trace-weight-lower-bound}
             w>w(19/5)=\frac{153861525}{421654016}>\frac13.
\end{equation}

Set
\[
 Z_\pm=(1+w)W_\pm-1,\qquad G=Z_+Z_-.
\]
Both $Z_\pm\ge w>0$, and expansion shows that $V^*=1$ is
equivalent to $G=1$. For each product define its logarithmic
derivative with respect to $\alpha$ by
\[
 P_\pm=\sum_k\frac{\alpha c_{\pm,k}}{1+\alpha c_{\pm,k}}.
\]
The elementary product inequality gives
\begin{equation}\label{eq:unit-trace-elasticity}
                         0\le P_\pm\le W_\pm-1.
\end{equation}

At a point with $G=1$, put $S=2+Z_++Z_->0$. The reciprocal
relations $Z_+^{-1}=Z_-$ and $Z_-^{-1}=Z_+$ give
\begin{equation}\label{eq:unit-trace-log-derivative}
 \frac{G'}G=\frac{w'}{1+w}S+
 \frac{\alpha'}\alpha
              \sum_\pm P_\pm(1+Z_\pm^{-1}).
\end{equation}
If $\alpha'\le0$, the first term is strictly negative and the
second is nonpositive, so $G'<0$.

If $\alpha'>0$, use~\eqref{eq:unit-trace-elasticity} and
$W_\pm=(Z_\pm+1)/(1+w)$ to obtain
\[
 \sum_\pm P_\pm(1+Z_\pm^{-1})
 \le \frac{1-w}{1+w}S.
\]
Here $w\le1$, because $Z_\pm\ge w$ and $Z_+Z_-=1$.
Equation~\eqref{eq:unit-trace-log-derivative} consequently implies
\[
 \frac{G'}G\le\frac{S}{1+w}
       \left(w'+(1-w)\frac{\alpha'}\alpha\right).
\]
The expression in parentheses equals
\[
                 (2w-1)f+\frac{4-12w}{\delta}.
\]
By~\eqref{eq:unit-trace-weight-lower-bound}, $w>1/3$.
For $1/3<w<1/2$ both terms are negative. For $w\ge1/2$,
the bound $f<6/\delta$ gives a value at most $-2/\delta<0$.
Thus every root of $G-1$ is a strict downward crossing.

As $t\to1^+$, $G\ge w^2\to\infty$. As $t\to\infty$,
$w,\alpha\to0$, each $W_\pm\to1$, and $G\to0$.
Continuity supplies a root; strict downward crossing at every
root gives uniqueness. The final assertion follows because $t$
and $at$ are distinct points of the same interval.
\end{proof}

This obstruction applies to the specialization requiring unit trace
at both conditional detector arguments when both exceed one.
It imposes no restriction on other real argument intervals or on
the full prescribed-trace norm equations with a varying rational
ordinate. The theorem asserts existence of a real root only;
its rationality is not asserted.

\subsection{A uniform scale bound for actual conditional tuples}

Retain the six-contact functions and notation of
Section~\ref{sec:unit-trace-positive-interval}. At an actual conditional
point tuple, the products contain the two tests
$f(H)=1/x(H)$ and $\overline{f(H-P)}$, with their conjugate-slot
counterparts, where
\[
 P=\left(\frac{57-28\sqrt2}{8},
          \frac{-660+417\sqrt2}{32}\right),\qquad
 E:\ y^2=x^3-(3-\sqrt2)^2x.
\]
All point and difference-chart guards are retained.

\begin{lemma}\label{lem:reference-real-separation}
For every guarded real point $H$, at the embedding $\sqrt2>0$,
\[
 \max\{|f(H)|,|f(H-P)|\}>\frac18.
\]
At the conjugate embedding the corresponding lower bound is $1/48$.
\end{lemma}
\begin{proof}
Write $P=(p,q)$ and the curve as $y^2=x^3-Dx$. For $H=(x,y)$
with $x\ne p$, the addition formula gives
\[
 x(H-P)=
 \frac{px^2+(p^2-D)x-Dp+2qy}{(x-p)^2}.
\]
For $x>0$ one has $|y|\le x^{3/2}$. At the first embedding,
\[
 0<p<\frac{11}{5},\quad |q|<\frac94,\quad
 p^2-D<\frac94,\quad \sqrt D<2.
\]
For $x\ge8$, using $\sqrt8>14/5$, the displayed formula yields
\[
 x(H-P)<\frac{1600}{841}
       \left(\frac{11}{5}+\frac9{32}+\frac{45}{28}\right)
       =\frac{45790}{5887}<8.
\]
Every negative real abscissa belongs to $[-\sqrt D,0]$.
The same negative-component bound applies to $H-P$. Thus
$|x(H)|\ge8$ forces $|x(H-P)|<8$, proving the first assertion.

At the conjugate embedding,
\[
 0<p<\frac{49}{4},\quad |q|<\frac{157}{4},\quad
 p^2-D<127,\quad\sqrt D<5.
\]
For $x\ge48$, use $\sqrt{48}>69/10$ to obtain
\[
 x(H-P)<\left(\frac{192}{143}\right)^2
   \left(\frac{49}{4}+\frac{127}{48}+\frac{785}{69}\right)
   =\frac{22275840}{470327}<48.
\]
The negative component has $|x|<5$, which proves the second assertion.
\end{proof}

\begin{proposition}\label{prop:unit-trace-positive-scale}
At an actual conditional point tuple, if $a>1$, $t>0$, and
$V^*(t)=V^*(at)=1$, then $a<147456$.
\end{proposition}
\begin{proof}
Theorem~\ref{thm:unit-trace-positive-unique} excludes $t>1$,
and $t=1$ is outside the source domain. Suppose $0<t<1$.
At a unit-trace point, $Z_+Z_-=1$ and $Z_\pm\ge w$ imply
$W_\pm\le1/w$. In particular $w<1$, since the actual reciprocal
tests are nonzero.

One of $W_+^{-1},W_-^{-1}$ is at least $(1+w)/2$, so the corresponding
product $W_{\rm good}$ is below two. Every test in that slot has
absolute value less than $\alpha^{-1/2}$. If $\alpha\ge2304$,
Lemma~\ref{lem:reference-real-separation}, applied at the two embeddings,
forces the complementary slot to contain two test squares exceeding
$1/64$ and $1/2304$, respectively. Hence
\[
 W_{\rm bad}>
  (1+\alpha/64)(1+\alpha/2304)>
                         \frac{\alpha^2}{147456}.
\]
The equation $W_+^{-1}+W_-^{-1}=1+w$ gives
$w\le W_{\rm bad}^{-1}$. Since $w\alpha=1/(4(t-1)^4)$,
\[
 w<\frac{147456}{\alpha^2}
   =2359296w^2(t-1)^8,
 \qquad
                  w>\frac1{9\cdot2^{18}}.
\]
If $\alpha<2304$, the same final bound follows from
$w=1/(4\alpha(t-1)^4)>1/9216$.

The function $w(t)/t$ is increasing on $(0,1)$, and
\[
 \left.\frac{w(t)}t\right|_{t=1/16}
       =\frac{12237824}{854296875}<\frac1{64}.
\]
If $t<1/16$, it follows that $t>64w>1/(9\cdot2^{12})$;
the same final bound is immediate if $t\ge1/16$.

Every test square at this inner root also satisfies
\[
 c_{\pm,k}\le\frac{1/w-1}{\alpha}
           =4(t-1)^4(1-w)<4.
\]
For $s\ge4$, both $w(s)$ and $\alpha(s)$ decrease, and
\[
 w(4)=\frac{2431}{8748},\qquad
 \alpha(4)=\frac{27}{2431}.
\]
As the native product has six factors,
\[
 (1+w(s))W_\pm(s)
 \le\frac{11179}{8748}\left(\frac{2539}{2431}\right)^6<2.
\]
Thus $V^*(s)>1$ for $s\ge4$, and the second unit-trace root obeys
$at<4$. Combining this with $t>1/(9\cdot2^{12})$ proves
$a<9\cdot2^{14}=147456$.
\end{proof}

\begin{lemma}\label{lem:unit-trace-negative-outer}
For an actual conditional point tuple, $V^*(t)=1$ has exactly one
real root on $t<-1$.
\end{lemma}
\begin{proof}
Put $t=-x$, $x>1$, $d=x+1$, and
\[
 F_-(x)=F(-x)=\prod_{c\in\{0,1,1/2,1/3,1/4,2/5\}}(x-c).
\]
Now $w=F_-/(4d^8)$ and $\alpha=d^4/F_-$. Derivatives in this
proof are with respect to $x$. Write $f=F_-'/F_-$ and
$g=\alpha'/\alpha=4/d-f$. Then $f>6/d$, $g<0$, and $0<w<1/4$.

At $G=Z_+Z_-=1$, put $S=2+Z_++Z_-$ and retain $P_\pm$ from
the positive-interval proof. The product inequality
$P_\pm\ge1-W_\pm^{-1}$ gives
\[
 \sum_\pm P_\pm(1+Z_\pm^{-1})\ge2+2w-wS.
\]
Since $g<0$,
\[
 \frac{G'}G\le
 \left(\frac{w'}{1+w}-wg\right)S+2(1+w)g.
\]
The coefficient of $S$ is positive, because
\[
 \frac{w'}{1+w}-wg
 =\frac{w}{1+w}
   \left((2+w)f-\frac{12+4w}{d}\right)
 >\frac{2w^2}{(1+w)d}>0.
\]
The bounds $w\le Z_\pm\le1/w$ give $S\le(1+w)^2/w$.
Consequently
\[
 \frac{G'}G\le(1+w)\left(wf-\frac{4(1+w)}d\right)<0.
\]
For the last inequality, $F_-'<6x^5$ implies
$wf<3x^5/(2d^8)<4/d$.

Each actual slot contains the nonzero test $f(H)$, so $G\to\infty$
as $x\to1^+$. At infinity $w,\alpha\to0$ and $G\to0$.
Every root is a strict downward crossing, proving existence and uniqueness.
\end{proof}

\begin{proposition}\label{prop:unit-trace-negative-scale}
At an actual conditional point tuple with positive weights, if
$a>1$, $t<0$, and $V^*(t)=V^*(at)=1$, then $a<80$.
\end{proposition}
\begin{proof}
Lemma~\ref{lem:unit-trace-negative-outer} excludes $t<-1$.
Positivity of $F(t)$ leaves only
\[
          -\frac12<t<-\frac25\qquad\hbox{or}\qquad
          -\frac13<t<-\frac14.
\]
Thus $1/4<|t|<1/2$ and $|t-1|<3/2$. At the first unit-trace root,
every test square is below $4(t-1)^4<81/4$. For $x\ge20$,
\[
 w(-x)<\frac1{4x^2}\le\frac1{1600},\qquad
 \alpha(-x)<\frac{(x+1)^4}{(x-1)^6}
                         \le\frac{21^4}{19^6}.
\]
It follows that
\[
 (1+w(-x))W_\pm(-x)
 <\frac{1601}{1600}
       \left(1+\frac{81}{4}\frac{21^4}{19^6}\right)^6<2.
\]
Hence $V^*(-x)>1$ for $x\ge20$. The second unit-trace root has
$|at|<20$, giving $a<20/|t|<80$.
\end{proof}

\begin{corollary}\label{cor:unit-trace-conditional-scale}
At every actual conditional tuple with prepared positive weights,
simultaneous normalized unit trace at arguments $t,at$, $a>1$,
requires $a<147456$, regardless of the real slope interval.
For either conditional source, with
\[
       a=(X^2+B)^d+\Tr(x(A))^d,\qquad d\in\{4,6\},
\]
this specialization is therefore impossible for every nonzero integer
input when $B\ge19$.
\end{corollary}
\begin{proof}
Combine Propositions~\ref{prop:unit-trace-positive-scale}
and~\ref{prop:unit-trace-negative-scale}. For a nonzero integer $X$
and $B\ge19$, $a\ge20^4=160000>147456$.
\end{proof}

The corollary excludes the specialization requiring both normalized
traces to equal one. It does not exclude the full adaptive norm
equations, whose rational ordinates remain free within their finite
algebraic adjunctions.

\section{Exact checks and remaining scope}

The companion verifier
\begin{center}\small
\nolinkurl{negative_results/computation/fixed_traces_v51/check.py}
\end{center}
uses only the standard library and checks primality, the exact
scalar contact valuations, the disjoint residue sets, sample exponents
including $20483$, the native transformation identities in rational
pair arithmetic, and the scalar-prepared norm-orbit examples. The
universal quantifiers over exponents, slopes, auxiliary tuples and
normalization sequences follow from the proofs above, not from the
finite computations. The separate script
\begin{center}\small
\nolinkurl{negative_results/computation/fixed_traces_v51/check_chamber.py}
\end{center}
uses exact symbolic arithmetic to check the derivative identities, the
positive shifted polynomial, and the rational weight bound in the
conditional interval theorem. The scripts \texttt{check\_positive\_scale.py}
and \texttt{check\_negative\_scale.py} in the same directory check the
exact reference-point bounds and the remaining scale inequalities.

The flexible prescribed-trace equations
\[
                    (m_i\mathcal V_i)^2-d_i y_i^2=4\omega_i/17
                         \qquad(i=0,1)
\]
from~\cite{Adaptive} are not restricted to constant traces. Their
simultaneous rational realization on actual detector tuples remains
open. The present obstructions neither settle that assertion nor
establish a lower bound on the number of variables in a polynomial
defining the integers over the rationals.

\begin{thebibliography}{9}\raggedright
\bibitem{Unit}
\emph{A single norm determinant for the 25- and 13-variable constructions},
repository source,
\nolinkurl{papers/finite_algebra/unit_norm_20261114_v28.tex}.
\bibitem{Adaptive}
\emph{Adaptive levels and varying norm classes for native elliptic detectors},
repository source,
\nolinkurl{papers/pole_parameters/adaptive_native_complements_20261116_v50.tex}.
\bibitem{Finite}
\emph{A finite-class obstruction for the native scalar source},
repository source,
\nolinkurl{negative_results/papers/pole_parameters/finite_nonsquare_class_obstruction_20261116_v47.tex}.
\bibitem{NormScalar}
\emph{Norm preparation for general and conditional elliptic detectors},
repository source,
\nolinkurl{papers/pole_parameters/norm_compatible_scalar_20261116_v47.tex}.
\bibitem{Conditional}
\emph{Thirteen Variables from Direct Elliptic Index Comparison},
repository source,
\nolinkurl{papers/conditional_few_variables/direct_input13_20261110_v22.tex}.
\end{thebibliography}
\end{document}
